\section{Sampling Theory}
\label{sec:sampling}

\subsection{Latent-Family Variances}

Let $\hat\theta$ be the maximum-likelihood estimator of $\theta = \log u$ from $n$ draws of the latent family \cref{eq:geo_family}.
This estimator exists only for $u\in(0,1)$.
Standard exponential-family theory gives $\sqrt{n}(\hat\theta-\theta)\xrightarrow{d}\mathcal{N}(0,1/\mathcal{I}(\theta))$ with Fisher information $\mathcal{I}(\theta) = \eta(\eta-1)$.
Fisher information is non-negative on this family.
It is not defined for $\rho\le 0$, and it does not change sign.

Define $\psi = 2\arctanh\sqrt{u}$ for $u\in(0,1)$.
Then $\mathrm{d}\psi/\mathrm{d}\theta = \sqrt{\mathcal{I}(\theta)}$, so the delta method yields
\begin{equation}
  \Var(\hat\psi) \approx \frac{1}{n}
  \label{eq:psi_clt}
\end{equation}
for this latent maximum-likelihood estimator.
The same delta method uses $\mathrm{d}\eta/\mathrm{d}\theta = \eta(\eta-1)$, and therefore
\begin{equation}
  \Var(\hat\eta)
  \approx \frac{[\eta(\eta-1)]^{2}}{n\,\eta(\eta-1)}
  = \frac{\eta(\eta-1)}{n}.
  \label{eq:var_eta_latent}
\end{equation}
The variance is small when $\eta$ is near $1$ and grows as $\eta\to\infty$.
The earlier factor $\eta^{2}(\eta-1)/n$ is too large by $\eta$.

The coordinate $\psi$ is not Fisher's $z$.
At $\tau = 0$ one has $u = \rho$, so $\psi = 2\arctanh\sqrt{\rho}$ for $\rho\in(0,1)$, whereas $z = \arctanh\rho$.
The two transforms differ.

For display when $\rho<0$, the signed extension $\operatorname{sign}(\rho)\cdot 2\arctanh\sqrt{|u|}$ is continuous at $\rho = 0$.
Its derivative in $u$ diverges as $u\to 0$.
The extension is not the arc length of \cref{eq:geo_family}, because that family does not exist on $\rho\le 0$.
There is then no maximum-likelihood estimator in the latent family and no $1/n$ law.

\subsection{Plug-in Variance under Bivariate Normality}

The estimator used on paired observations is not the latent-family maximum-likelihood estimator.
It is the plug-in $(\hat\kappa,\hat\rho)$ computed from $n$ pairs.
The next proposition is for that estimator.
\Cref{sec:numerical} compares it with a bivariate-normal simulation.

\begin{proposition}[Plug-in variance]
\label{prop:plugin}
Let $(X_{\mathrm{A}},X_{\mathrm{B}})$ be bivariate normal, and let $\hat\rho$ and $\hat\tau = \tfrac{1}{2}\log(s_{\mathrm{B}}^{2}/s_{\mathrm{A}}^{2})$ be computed from $n$ pairs.
Set $u = \rho\,\sech\tau$.
Asymptotically, $\Var(\hat\rho)\approx(1-\rho^{2})^{2}/n$, $\Var(\hat\tau)\approx(1-\rho^{2})/n$, and $\Cov(\hat\rho,\hat\tau) = o(1/n)$.
Therefore
\begin{equation}
  \Var(\hat u)
  \approx \frac{(\sech^{2}\tau - u^{2})(1-u^{2})}{n}.
  \label{eq:var_u_plugin}
\end{equation}
\end{proposition}

\begin{proof}
The argument is the delta method recorded in \cref{sec:appendix_plugin}.
The leading covariance of $\hat\rho$ and $\hat\tau-\tau$ vanishes because the swap $\mathrm{A}\leftrightarrow\mathrm{B}$ fixes $\hat\rho$ and reverses the sign of $\hat\tau-\tau$.
The variance of a log sample variance is $2/n$, and $\Cov(\log s_{\mathrm{A}}^{2},\log s_{\mathrm{B}}^{2})\approx 2\rho^{2}/n$, which gives $\Var(\hat\tau)$.
The partial derivatives are $\partial u/\partial\rho = \sech\tau$ and $\partial u/\partial\tau = -u\tanh\tau$.
Their combination simplifies by $\sech^{2}\tau\,(1-\rho^{2}) = \sech^{2}\tau - u^{2}$.
\end{proof}

\begin{corollary}
\label{cor:plugin_eta_psi}
On $u\in(0,1)$, the same delta method gives
\begin{align}
  \Var(\hat\eta)
  &\approx \frac{\eta^{4}\,(\sech^{2}\tau-u^{2})(1-u^{2})}{n},
  \label{eq:var_eta_plugin}\\
  \Var(\hat\psi)
  &\approx \frac{(\sech^{2}\tau-u^{2})(1+u)}{n\, u\,(1-u)}.
  \label{eq:var_psi_plugin}
\end{align}
\end{corollary}

\Cref{eq:var_u_plugin} depends on $\tau$, not on $u$ alone.
No transformation of $u$ can stabilise the plug-in variance exactly.
At $\tau = 0$ the formula reduces to $\Var(\hat\rho)\approx(1-\rho^{2})^{2}/n$, whose stabilising transform is Fisher's $z$.
\Cref{eq:var_psi_plugin} diverges as $u\to 0$.
The plug-in $\hat\psi$ is therefore least stable near independence.
That is where many of the sports indicators in \textcite{brownPEFempirical} sit.
The uniform law $\Var(\hat\psi)\approx 1/n$ is \cref{eq:psi_clt}.
It is not a property of $\hat\psi$ computed from $(\hat\kappa,\hat\rho)$.
